\scanpage{33}
on a mieux : % L1
\[
M\in\mathrm{Ob}\,\mathcal M^{+i}(X)\Longleftrightarrow M(-j)\in\mathrm{Ob}\,\mathcal M^{+i+2j}(X) % L2
\]

De cette façon, la filtration \inserted{de $\mathcal M^+(X)$} par les % L3
$\mathcal M^{+i}(X)$ \inserted{$(i\geq0)$} peut se prolonger en une filtration % L4
de $\mathcal M(X)$ par des $\mathcal M^i(X)$ $(i\in\mathbb Z)$ de % L5
telle façon que pour $i\geq0$, on \unsure{suit|sent|soit} la filtration % L6
déjà \deleted{\unsure{$\mathcal M^+(X)$|$\mathcal M(X)$}} envisagée, et pour $i$ quelconque, % L7
on définit % L8
\[
M\in\mathrm{Ob}\,\mathcal M^i(X)\Longleftrightarrow M(-j)\in\mathrm{Ob}\,\mathcal M^{i+2j}(X) % L9
\]
où on prend $j$ assez grand pour que % L10
$M(-j)\in\mathrm{Ob}\,M^+(X)$. On notera que % L11
\[
\mathcal M^{+i}(X)=\mathcal M^i(X)\cap\mathcal M^+(X) % L12
\]
% slanted
\margin{indices en bas} % L12a
si $i\geq0$, et pour tt $i$ si on définit % L13
$\mathcal M^{\unsure{+i|-i}}(X)=\{0\}$\deleted{\unsure{$\mathcal M^+(X)$}} si $i<0$ \deleted{(\unsure{ce qui revient au|ce --- pour les} % L14
même si on veut que $\mathcal M^{+i}(X)$ \inserted{$\subset\mathcal M^+(X)$ et} \unsure{existe aussi ici|consiste aussi ici}}] % L15

On aura % L16
\[
\begin{aligned}
\otimes\deleted{\mathbb Z}\mathbb Q(-1)&:\quad\mathcal M^i(X)\overset{\text{équiv}}{\simeq}\mathcal M^{i+2}(X)\\ % L17
\otimes\deleted{\mathbb Z}\mathbb Q(j)&:\quad\mathcal M^i(X)\overset{\text{équiv}}{\simeq}\mathcal M^{\unsure{i-2j|i+2j}}(X) % L18
\end{aligned}
\]

% archive sheet number: 32
